\section{Introduction}
Answering a geographic question often requires more than retrieving a passage that resembles the query. A request such as finding amenities near a named place can require resolving the place to one or more observed geometries, retrieving candidate features, checking category evidence, carrying out a spatial operation, and combining the results without losing provenance. These operations depend on different data and tools. If an anchor cannot be resolved, a distance calculation is not justified; if a category is absent from a source, proximity alone does not establish category membership. A system that hides these dependencies inside a single prompt makes it difficult to inspect why an answer was produced or where an investigation failed.

Existing geospatial question-answering work has exposed this gap. MapQA pairs questions with geospatial entities and geometries, and reports that retrieval systems can capture some semantic and spatial cues while explicit calculations and multi-hop reasoning remain challenging~\cite{li2025mapqa}. Spatial-RAG combines spatial-database retrieval with semantic retrieval, but its published evaluation is centered on a tourism setting~\cite{yu2025spatialrag}. These results motivate systems that treat spatial operations, source coverage, and semantic matching as distinct evidence-producing steps rather than expecting a language model to infer them from text alone.

We present \emph{Agentic Tree-RAG}, a framework for planning and executing evidence-grounded investigations over heterogeneous geographic data. A planner constructs a query-specific tree whose nodes represent retrieval, computation, or synthesis tasks. Parent-child edges encode execution control, while explicit input references encode evidence dependencies. The plan supports both unconditional work that remains useful across outcomes and conditional branches activated only when their prerequisites hold. NodeAgents retrieve source candidates; a MainAgent invokes registered tools; JEV makes evidence-conditioned tree-operation decisions; and a single QueryTreeManager validates and applies state transitions. A separate Evidence Tree preserves retrieved source pointers and tool outputs, including evidence from branches later pruned from the active plan. On the data side, a Description Tree organizes source descriptions and pointers, with BM25 and modality-specific HNSW retrieval combined by reciprocal-rank fusion and reranking.

The framework is designed around three requirements. First, a plan must expose real data dependencies: a spatial computation consumes the observations needed to perform it, rather than relying on undeclared context. Second, conditional execution must preserve unresolved and alternative outcomes instead of silently treating missing evidence as a negative fact. Third, every answer claim should be traceable to retrieved source records or tool results. For geospatial data, the system can use GeoSPARQL when a source has an executable, verified ontology and geometry binding; ontology semantics and spatial functions supplement source retrieval, while category membership remains grounded in source evidence~\cite{ogc2023geosparql}.

This work makes two contributions:
\begin{itemize}
  \item We formulate geospatial question answering as a source-grounded investigation with explicit control flow and evidence dependencies, separating the execution plan from the data-source description index.
  \item We describe an agent architecture that combines conditional query-tree execution, asynchronous branch scheduling, hybrid heterogeneous retrieval, tool-mediated spatial and tabular operations, source-grounded ontology alignment, and a persistent evidence structure.
\end{itemize}

The remainder of the paper is organized as follows. Section~\ref{sec:related-work} reviews geospatial QA, spatial retrieval, and tool-using agents. Section~\ref{sec:method} presents the query-tree representation, execution semantics, source retrieval, ontology integration, spatial tools, and evidence-grounded synthesis.
